
\documentclass{article}
\usepackage{colortbl}
\usepackage{makecell}
\usepackage{multirow}
\usepackage{supertabular}

\begin{document}

\newcounter{utterance}

\centering \large Interaction Transcript for game `function\_detective', experiment `list\_passive\_labeled\_pairs\_oneshot', episode 107 with mistral{-}large{-}3{-}azure.
\vspace{24pt}

{ \footnotesize  \setcounter{utterance}{1}
\setlength{\tabcolsep}{0pt}
\begin{supertabular}{c@{$\;$}|p{.15\linewidth}@{}p{.15\linewidth}p{.15\linewidth}p{.15\linewidth}p{.15\linewidth}p{.15\linewidth}}
    \# & \multicolumn{2}{c}{Player} && \multicolumn{2}{c}{Game Master} \\
    \hline

    \theutterance \stepcounter{utterance}  
    & & & \multicolumn{4}{p{0.6\linewidth}}{
        \cellcolor[rgb]{0.9,0.9,0.9}{
            \makecell[{{p{\linewidth}}}]{
                \texttt{\tiny{[P1$\langle$GM]}}
                \texttt{I have a hidden Python function for you to figure out.} \\
\\ 
\texttt{Your job is to infer what it does and then submit a correct Python implementation.} \\
\\ 
\texttt{Function signature:} \\
\texttt{(items: List{[}int{]}) {-}> int} \\
\\ 
\texttt{You have 1 turns total.} \\
\\ 
\texttt{I have already given you a fixed list of candidate input{-}output pairs with true or false labels. Solve using those statements only.} \\
\\ 
\texttt{You must make your best SOLVE attempt no later than the second{-}to{-}last turn.} \\
\texttt{Do not wait until the very end.} \\
\\ 
\texttt{You can use these actions:} \\
\texttt{{-} JSON format described at the end of the prompt} \\
\\ 
\\ 
\texttt{Do not add import statements in your guess.} \\
\texttt{Anything you might need is already available, so just write the function body in hypothesis and return the answer.} \\
\\ 
\texttt{Preloaded examples:} \\
\texttt{1. (({[}2, 1{]}, 2), True)} \\
\texttt{2. (({[}1, 2, 3{]}, 9), False)} \\
\texttt{3. (({[}0{]}, 0), True)} \\
\texttt{4. (({[}5, {-}5, 5{]}, {-}49), False)} \\
\texttt{5. (({[}{-}10{]}, 0), True)} \\
\texttt{6. (({[}{]}, {-}1), False)} \\
\texttt{7. (({[}1, 1{]}, 1), True)} \\
\texttt{8. (({[}3, 2, 1{]}, 9), False)} \\
\texttt{9. (({[}{-}1, 0, 1{]}, 0), True)} \\
\texttt{10. (({[}{-}2, 0, 0{]}, {-}1), False)} \\
\\ 
\texttt{Scoring:} \\
\texttt{{-} correctness is required for any non{-}zero score} \\
\texttt{{-} in interactive modes, solving with fewer information requests is better} \\
\\ 
\texttt{Input formatting reminders:} \\
\texttt{{-} strings must be quoted} \\
\texttt{{-} booleans must be True/False} \\
\texttt{{-} use commas between multiple values} \\
\texttt{{-} do not write any extra action tags in the same message} \\
\\ 
\texttt{Format example:} \\
\texttt{You write:} \\
\texttt{\{'current\_mode': <probing/confirming/don't know; choose whichever best describes your current state in the task>, 'input\_rationale': <brief description of why you chose this hypothesis>, 'hypothesis': <either the exact string "unknown" if you have no guess yet OR your current best guess as runnable Python matching the given signature, with a one{-}line docstring stating what it does. Never write placeholder code, if you cannot commit to a guess, use "unknown". Match the given signature.} \\
\texttt{e.g.} \\
\texttt{def f(x):} \\
\texttt{	"""Returns the absolute value of x."""} \\
\texttt{	return x if x >= 0 else {-}x} \\
\\ 
\texttt{Use "unknown" if you have no guess.>, 'confidence\_score': <score between 1{-}100. 100 means highest confidence in your hypothesis>\}} \\
\\ 
\texttt{Output only one valid json per turn.} \\
            }
        }
    }
     \\ \\

    \theutterance \stepcounter{utterance}  
    & \multicolumn{4}{p{0.6\linewidth}}{
        \cellcolor[rgb]{0.561,0.737,0.561}{
            \makecell[{{p{\linewidth}}}]{
                \texttt{\tiny{[P1$\rangle$GM]}}
                \texttt{```json} \\
\texttt{\{} \\
\texttt{    "current\_mode": "probing",} \\
\texttt{    "input\_rationale": "Analyzing the input{-}output pairs, the function seems to return the square of the sum of the list if the list has an odd number of elements, but the examples show some exceptions. For instance, ({[}2, 1{]}, 2) returns True (sum is 3, square is 9, but output is 2), which doesn't fit. Alternatively, it might return the product of the elements if the list length is odd, but ({[}2, 1{]}, 2) would be 2*1=2 (matches), ({[}{-}10{]}, 0) would be {-}10 (but output is 0, so no). Another pattern: the function might return the sum of the squares of the elements if the list length is odd, but ({[}2, 1{]}, 2) would be 4+1=5 (doesn't match). The examples ({[}{-}10{]}, 0) and ({[}0{]}, 0) suggest that if the only element is negative or zero, the output is 0. The example ({[}1, 1{]}, 1) suggests the product of the list if the length is even, but ({[}2, 1{]}, 2) is length 2 and product is 2 (matches), but ({[}3, 2, 1{]}, 9) is length 3 and product is} \\
            }
        }
    }
    & & \\ \\

    \theutterance \stepcounter{utterance}  
    & & & \multicolumn{2}{p{0.3\linewidth}}{
        \cellcolor[rgb]{0.9,0.9,0.9}{
            \makecell[{{p{\linewidth}}}]{
                \texttt{\tiny{[GM$|$GM]}}
                \texttt{Game Verdict: LOSS (turn limit reached)} \\
            }
        }
    }
    & & \\ \\

\end{supertabular}
}

\end{document}
